\section{Relative entropies and binary hypothesis testing}
\label{sec:hypothesis_testing}

This section compares the relative entropy with the quotients built from the
ordinary overlap, determines the optimal error for discriminating two states
by one measurement, and bounds the trace distance below by the overlap.

\begin{definition}[Tsallis relative entropy]
    \label{def:tsallis_relative_entropy}
    \lean{Matrix.tsallisRelativeEntropy}
    \leanok
    For positive semidefinite $\rho,\sigma\in\MN d$ and real $\alpha\ne1$, set
    \begin{align}
        S^{(\alpha)}(\rho|\sigma)
            &=\frac{\operatorname{Re}\tr[\rho-\rho^\alpha\sigma^{1-\alpha}]}{1-\alpha}.
        \label{eq:tsallis_relative_entropy}
    \end{align}
    Powers are defined by spectral calculus; a negative power inverts on the
    support and vanishes on the kernel, and $X^0=\Id$.
\end{definition}

\begin{lemma}[Spectral double sums]
    \label{lem:trace_function_double_sum}
    \lean{Matrix.re_trace_cfc_mul_cfc_eq_double_sum,
        Matrix.eigenvalues_mul_normSq_eq_zero_of_ker_le}
    \leanok
    Let $\rho=\sum_i\lambda_i\ketbra{x_i}{x_i}$ and
    $\sigma=\sum_j\mu_j\ketbra{y_j}{y_j}$ be Hermitian, with orthonormal
    eigenbases, and put $P_{ij}=\lvert\braket{x_i}{y_j}\rvert^2$. For real
    functions $f,g$,
    \begin{align}
        \operatorname{Re}\tr[f(\rho)g(\sigma)]
            &=\sum_{i,j}f(\lambda_i)\,g(\mu_j)\,P_{ij}.
        \label{eq:trace_function_double_sum}
    \end{align}
    If $\rho\ge0$ and $\ker\sigma\subseteq\ker\rho$, then $\lambda_iP_{ij}=0$
    whenever $\mu_j=0$.
    This is \cite[Equations~(8.36)--(8.37)]{Wolf2012Quantum}.
\end{lemma}
\begin{proof}
    \leanok
    Conjugating by the two eigenvector unitaries reduces the trace to a
    product of two diagonal matrices and the overlap matrix
    $W_{ij}=\braket{x_i}{y_j}$, which gives
    (\ref{eq:trace_function_double_sum}). For the support statement apply
    (\ref{eq:trace_function_double_sum}) with $f(x)=x$ and $g$ the indicator
    of $\{0\}$. The trace equals the sum of the quadratic forms of $\rho$ on
    the eigenvectors of $\sigma$ with eigenvalue zero, which vanish by the
    kernel condition. All terms of the double sum are nonnegative, so each
    vanishes.
\end{proof}

\begin{theorem}[Relation among relative entropies]
    \label{thm:relative_entropy_relations}
    \lean{Matrix.sub_rpow_mul_rpow_le_mul_log,
        Matrix.re_trace_sub_le_mul_quantumRelativeEntropy,
        Matrix.tsallisRelativeEntropy_le_quantumRelativeEntropy,
        Matrix.quantumRelativeEntropy_le_tsallisRelativeEntropy}
    \leanok
    \uses{def:tsallis_relative_entropy, def:quantum_relative_entropy}
    Let $\rho,\sigma\in\MN d$ be positive semidefinite with
    $\ker\sigma\subseteq\ker\rho$. For every real $a$,
    \begin{align}
        \operatorname{Re}\tr[\rho-\rho^a\sigma^{1-a}]
            &\le(1-a)\,S(\rho\|\sigma).
        \label{eq:relative_entropy_relations_common}
    \end{align}
    Consequently, for $\alpha<1<\beta$,
    \begin{align}
        S^{(\alpha)}(\rho|\sigma)\le S(\rho\|\sigma)\le S^{(\beta)}(\rho|\sigma).
        \label{eq:relative_entropy_relations}
    \end{align}
    This is \cite[Equation~(8.38)]{Wolf2012Quantum} restricted to finite
    relative entropy and to exponents different from one. When the kernel
    condition fails the relative entropy is infinite; with the source's
    assignment of infinite values for $\alpha<0$ the left inequality fails.
    The comparison is recorded in the paper-gap notes
    \path{docs/paper-gaps/wolf_ch8_relative_entropy_relations.tex} and
    \path{docs/paper-gaps/wolf_ch8_tsallis_negative_exponent_footnote.tex}.
\end{theorem}
\begin{proof}
    \leanok
    \uses{lem:trace_function_double_sum}
    By Lemma~\ref{lem:trace_function_double_sum} every trace in
    (\ref{eq:relative_entropy_relations_common}) is a double sum against
    $P_{ij}$, and the terms with $\lambda_i>0$, $\mu_j=0$ have $P_{ij}=0$.
    For $\lambda,\mu>0$ the inequality $\log x\le x-1$ at
    $x=(\mu/\lambda)^{1-a}$ gives
    \begin{align}
        \lambda-\lambda^a\mu^{1-a}
            &\le(1-a)\,\lambda(\log\lambda-\log\mu),
        \label{eq:relative_entropy_relations_scalar}
    \end{align}
    which is \cite[Equation~(8.39)]{Wolf2012Quantum} with $s=1-a$ or
    $s=a-1$. For $\lambda=0$ the left side is $-0^a\mu^{1-a}\le0$ and the
    right side vanishes. Summing gives
    (\ref{eq:relative_entropy_relations_common}); dividing by $1-\alpha>0$ and by
    $1-\beta<0$ gives the two halves of (\ref{eq:relative_entropy_relations}).
\end{proof}

\begin{definition}[Error of a binary test]
    \label{def:hypothesis_test_error}
    \lean{Matrix.hypothesisTestError}
    \leanok
    For matrices $\rho,\sigma,P\in\MN d$ and a prior $p\in\mathbb R$, the
    average error of guessing $\rho$ on the effect $P$ and $\sigma$ on
    $\Id-P$ is
    \begin{align}
        e(P)&=p\operatorname{Re}\tr[\rho(\Id-P)]
            +(1-p)\operatorname{Re}\tr[\sigma P].
        \label{eq:hypothesis_test_error}
    \end{align}
    This is \cite[Equation~(8.40)]{Wolf2012Quantum}.
\end{definition}

\begin{theorem}[Quantum Neyman--Pearson]
    \label{thm:quantum_neyman_pearson}
    \lean{Matrix.hypothesisTestError_sub_eq, Matrix.hypothesisTestError_ge,
        Matrix.hypothesisTestError_ge_traceNorm,
        Matrix.hypothesisTestError_eq_iff,
        Matrix.hypothesisTestError_posPartSupportProj,
        Matrix.isLeast_hypothesisTestError}
    \leanok
    \uses{def:hypothesis_test_error, def:trace_norm}
    Let $\rho,\sigma\in\MN d$ be Hermitian with unit trace, let $p\in\mathbb R$,
    and let $A=p\rho-(1-p)\sigma=A_+-A_-$ be the decomposition into positive
    and negative parts. For every effect $0\le P\le\Id$,
    \begin{align}
        e(P)&\ge\tfrac12\bigl(1-\lVert A\rVert_1\bigr).
        \label{eq:quantum_neyman_pearson}
    \end{align}
    Equality holds if and only if $PA_+=A_+$ and $PA_-=0$. The projection onto
    the range of $A_+$ attains equality, so the right side of
    (\ref{eq:quantum_neyman_pearson}) is the least error over all effects.
    This is \cite[Equation~(8.41)]{Wolf2012Quantum}, for all density matrices
    and priors, with the equality clause corrected as in
    \path{docs/paper-gaps/wolf_ch8_neyman_pearson_kernel.tex}.
\end{theorem}
\begin{proof}
    \leanok
    With $\tr\rho=\tr\sigma=1$, a direct computation gives
    \begin{align}
        e(P)-\tfrac12\bigl(1-\lVert A\rVert_1\bigr)
            &=\operatorname{Re}\tr[(\Id-P)A_+]+\operatorname{Re}\tr[PA_-],
        \label{eq:quantum_neyman_pearson_excess}
    \end{align}
    using $\lVert A\rVert_1=\tr A_++\tr A_-$ and $\tr A=2p-1=\tr A_+-\tr A_-$.
    For an effect both terms are traces of products of positive semidefinite
    matrices, hence nonnegative. Such a trace vanishes only if the product
    vanishes, which gives the equality condition. Writing $\Pi_+$ for the
    projection onto the range of $A_+$,
    \begin{align}
        \Pi_+A_+&=A_+, & \Pi_+A_-&=0,
        \label{eq:quantum_neyman_pearson_attainment}
    \end{align}
    so both terms of (\ref{eq:quantum_neyman_pearson_excess}) vanish at
    $P=\Pi_+$.
\end{proof}

\begin{corollary}[Equality clause and the kernel of the difference]
    \label{cor:quantum_neyman_pearson_kernel}
    \lean{Matrix.hypothesisTestError_eq_iff_eq_posPartSupportProj,
        Matrix.hypothesisTestError_half_self}
    \leanok
    \uses{thm:quantum_neyman_pearson}
    In the setting of Theorem~\ref{thm:quantum_neyman_pearson}, if $A$ is
    invertible then an effect attains equality in
    (\ref{eq:quantum_neyman_pearson}) if and only if it is the projection onto
    the range of $A_+$. Without invertibility this fails: for $p=\tfrac12$ and
    $\rho=\sigma$ every effect attains equality.
\end{corollary}
\begin{proof}
    \leanok
    If $PA_+=A_+$ and $PA_-=0$, then $PA=A_+$, which is also the product of
    the projection onto the range of $A_+$ with $A$; cancel $A$. For $p=\tfrac12$
    and $\rho=\sigma$, $A=0$ and $e(P)=\tfrac12$ for every $P$.
\end{proof}

\begin{lemma}[Trace distance and the overlap]
    \label{lem:trace_norm_overlap_bound}
    \lean{Matrix.re_trace_sub_re_trace_rpow_mul_rpow_le_posPart,
        Matrix.re_trace_add_sub_two_mul_re_trace_rpow_mul_rpow_le,
        Matrix.re_trace_add_sub_two_mul_re_trace_rpow_mul_rpow_le_traceNorm}
    \leanok
    \uses{def:trace_norm}
    Let $A,B\in\MN d$ be positive semidefinite and $s\in[0,1]$. Then
    \begin{align}
        \operatorname{Re}\tr A-\operatorname{Re}\tr[A^sB^{1-s}]
            &\le\tr(A-B)_+,
        \label{eq:trace_norm_overlap_posPart}\\
        \lVert A-B\rVert_1
            &\ge\operatorname{Re}\tr[A+B]-2\operatorname{Re}\tr[A^sB^{1-s}].
        \label{eq:trace_norm_overlap_bound}
    \end{align}
    The second inequality is \cite[Equation~(8.54)]{Wolf2012Quantum}. Singular
    matrices and the endpoints $s=0,1$ are included.
\end{lemma}
\begin{proof}
    \leanok
    Put $C=(A-B)_+$ and $M=B+C$. Then $A\le M$ and $B\le M$, so operator
    monotonicity of $X\mapsto X^t$ for $t\in[0,1]$ gives $A^t\le M^t$ and
    $B^t\le M^t$. Since $\tr[XY]\ge0$ for positive semidefinite $X,Y$,
    \begin{align}
        \tr A=\tr[A^sA^{1-s}]&\le\tr[A^sM^{1-s}],
        \label{eq:trace_norm_overlap_first}\\
        \tr[A^s(M^{1-s}-B^{1-s})]&\le\tr[M^s(M^{1-s}-B^{1-s})]
            =\tr M-\tr[M^sB^{1-s}],
        \label{eq:trace_norm_overlap_second}\\
        \tr B=\tr[B^sB^{1-s}]&\le\tr[M^sB^{1-s}].
        \label{eq:trace_norm_overlap_third}
    \end{align}
    Adding these gives (\ref{eq:trace_norm_overlap_posPart}), because
    $\tr M-\tr B=\tr C$. Finally
    $\lVert A-B\rVert_1=\tr(A-B)_++\tr(A-B)_-$ and
    $\tr(A-B)=\tr(A-B)_+-\tr(A-B)_-$ give
    $\tr[A+B]-\lVert A-B\rVert_1=2\tr A-2\tr(A-B)_+$, so
    (\ref{eq:trace_norm_overlap_posPart}) implies
    (\ref{eq:trace_norm_overlap_bound}).
\end{proof}
